\begin{table}[tb]
\centering\fontsize{8pt}{10pt}\selectfont\setlength{\tabcolsep}{2pt}
\caption{全24条件の成否と論理入出力token．不明生成を０へ置換せず，未評価・未実行も残す．}
\begin{tabular}{lrrr}
課題・反復 & M & H & L \\ \hline
写真1 & 成功 128,796 & 成功 135,628 & 成功 135,628 \\
写真2 & 未達 152,155 & 成功 135,628 & 成功 135,628 \\
精算1 & 成功 241,963 & 成功 211,110 & 成功 211,110 \\
精算2 & 成功 275,944 & 成功 211,110 & 成功 211,110 \\
返金1 & 成功 138,701 & 成功 102,087 & 成功 102,087 \\
返金2 & 成功 109,383 & 成功 102,087 & 成功 102,087 \\
書出し1 & 未達 257,266 & \shortstack[r]{未達 358,638\\＋不明1} & \shortstack[r]{未達 348,419\\＋不明1} \\
書出し2 & 未達 90,721 & \shortstack[r]{未達 358,638\\＋不明1} & \shortstack[r]{未達 346,789\\＋不明1} \\
\hline\end{tabular}

\end{table}
採用基準を満たした比較は0/4であり，両対照を両反復で上回るという固定基準は未達だった．条件を選び直した改善や，新手法の成立として扱わない．
反復1の固定基準はMに対し未達，Hに対し未達であった．
反復2の固定基準はMに対し未達，Hに対し未達であった．
反復1のM/H/Lの成功数は3/4，3/4，3/4であった．
反復2のM/H/Lの成功数は2/4，3/4，3/4であった．
全24条件で426生成，既知4,602,713論理token，使用量不明4生成，1156 API送信を記録した．所有推論器の稼働は16,263.578秒，うち読込10.061秒，評価41.470秒であった．仕事ごとの準備・作業・後片付けの合計は16,192.106秒であり，これらは重なりを持つため加算しない．
補助的な資源区間は本稿に含めず，欠測を含む全費用差を既知分だけから推定しない．
